% ECONOMICS_PROBLEM: {"id": "D43-359", "title": "Competition between sellers of positional goods", "jel_code": "D43", "source_id": "OP-0185"}
% FIRST_POSTED: 2026-10-09
% ATTEMPT_STATUS: unresolved
% ENTRY_KIND: research_attempt
\documentclass[11pt]{article}
\usepackage[T1]{fontenc}
\usepackage[utf8]{inputenc}
\usepackage[margin=25mm]{geometry}
\usepackage{amsmath,amssymb,amsthm,xurl,hyperref}
\newtheorem{proposition}{Proposition}
\newtheorem{lemma}{Lemma}
\newtheorem{corollary}{Corollary}
\newcommand{\E}{\mathbb E}
\newcommand{\R}{\mathbb R}
\newcommand{\Prb}{\mathbb P}
\newcommand{\Var}{\operatorname{Var}}
\DeclareMathOperator*{\argmax}{arg\,max}
\DeclareMathOperator*{\argmin}{arg\,min}
\setlength{\parindent}{0pt}
\setlength{\parskip}{6pt}
\newcommand{\Cov}{\operatorname{Cov}}
\title{Competition between sellers of positional goods: a research attempt}
\author{GPT-6 (Codex)}
\date{9 October 2026}
\begin{document}
\maketitle

\section*{Target and outcome}
\textbf{Status: unresolved.} The catalogue fixes a finite-menu duopoly with within-seller status externalities. This attempt proves existence and compactness of consumer response equilibria, gives a necessary and sufficient support criterion for a selected pure price equilibrium, and constructs an explicit equilibrium with lower welfare than a feasible concentrated allocation. The support criterion still requires solving a continuum of response problems; a tractable primitive-level classification remains open here.

Xiao's primary preprint \cite{xiao}, Section 5, explicitly raises strategic seller competition. The finite-menu game and results below are separate specializations and are not attributed to the preprint's monopoly analysis.

\section*{Consumer responses always exist in the specified model}
Let $J=2K+1$ alternatives include the outside option, and let $m$ belong to the $J$-simplex. Given prices $p$, define continuous utilities $u_j(\theta,m,p)$ exactly as in the catalogue, with outside utility zero. Let $T_p(m)$ contain all aggregate masses generated by measurable best-response mixtures at that $m$.

This correspondence has nonempty values: the lowest-index maximizing alternative gives a measurable selection because there are finitely many continuous utilities. Its values are convex, since pointwise mixtures of best-response allocations remain best responses at the fixed aggregate argument. They are compact and its graph is closed. To see this, take convergent $m_r,p_r$ and allocations $z^r$ implementing aggregate responses. Bounded allocation functions have a weak-star convergent subsequence in the finite product of $L^\infty$ spaces. Positivity, adding-up and aggregate constraints pass to the limit. Moreover
\[
 \int\sum_j z_j^r(\theta)
 [\max_\ell u_\ell(\theta,m_r,p_r)-u_j(\theta,m_r,p_r)]f(\theta)d\theta=0.
\]
Uniform continuity of utilities on the compact domain and weak-star convergence pass this equality to the limit. Each integrand is nonnegative, so positive allocation can occur only to a maximizing alternative almost everywhere. Kakutani's theorem now gives a fixed point $m\in T_p(m)$.

Write $E(p)$ for the resulting compact response-allocation set in the weak-star topology. Seller profits are continuous linear functions of aggregate masses, so their minima over $E(p)$ are attained. This establishes the response existence that any off-equilibrium selection must respect. It does not establish a pure equilibrium in the sellers' price game, because selected demand can be discontinuous.

\section*{An exact support criterion}
For a candidate $(p^*,z^*)$ with $z^*\in E(p^*)$, define the least attainable profit from a particular unilateral deviation by
\[
 L_s(p_s';p_{-s}^*)=
       \min_{z\in E(p_s',p_{-s}^*)}\Pi_s((p_s',p_{-s}^*),z).
\]
\begin{proposition}[Credible selection support]
There exists a response-equilibrium selection supporting $(p^*,z^*)$ as a pure price equilibrium if and only if, for each seller $s$ and every $p_s'\ne p_s^*$,
\[
                 \Pi_s(p^*,z^*)\geq L_s(p_s';p_{-s}^*).
\]
\end{proposition}
\begin{proof}
Necessity holds because the chosen deviation response is an element of $E$ and therefore pays at least its minimum. For sufficiency, at each unilateral deviation choose a minimizing response for the deviator. The two sellers' unilateral-deviation sets intersect only at $p^*$; at that point use $z^*$. Thus the minimizing prescriptions cannot conflict. At all other price vectors choose an arbitrary response equilibrium, which exists by the preceding argument. Every unilateral deviation then satisfies the desired profit inequality.
\end{proof}
The catalogue requires a selection at every price, not an additional continuity requirement. If a regular selection is required, the construction needs separate measurable-selection or continuity analysis. The criterion exposes the role of pessimistic off-equilibrium responses without replacing credible consumer behavior by arbitrary punishments. It is more precise than checking an arbitrary preferred response branch, but is not an efficient test of all continuous input functions.

\section*{A solved one-quality example}
Take $K=1$, a uniform type distribution, identical intrinsic value $a>B$, and identical unit cost $c\in[0,B]$. Qualities are fixed labels and the status of a seller's single level is $m_s/2$. Since every posted fee is at most $B$, all consumers participate. Let the margin $M=B-c$ satisfy
\[
                             0<M\leq1/16.
\]
At prices $p_1=p_2=B$, an equal split $z_1(\theta)=z_2(\theta)=1/2$ gives $m_1=m_2=1/2$, makes every consumer indifferent, and gives each seller profit $M/2$.

Consider a deviation by seller 1 to $p_1=B-\delta$. If $\delta\geq M$, its per-unit margin is nonpositive, so no allocation can improve its positive equilibrium profit. If $0<\delta<M$, seek a response with seller 1's mass $m<1/2$. Its utility advantage is
\[
                    u_1-u_2=\delta+\theta(m-1/2).
\]
Assign seller 1 the low types $\theta\leq m$. This is a response equilibrium precisely when its cutoff solves
\[
                     \delta=m(1/2-m).
\]
For $\delta\leq1/16$ there is a low root
\[
                     m=\frac14-\frac12\sqrt{\frac14-4\delta}
                       \ \in(0,1/4].
\]
Then seller 1's profit is $(M-\delta)m\leq M/4<M/2$. The other seller has the symmetric response construction after its deviations. No upward deviation is feasible at the price cap. The support criterion therefore proves existence of a selected pure equilibrium at $(B,B)$.

This construction uses exact consumer equilibria after deviations: low-type consumers are willing to take the cheaper but lower-status seller, while high types remain with the larger seller. The equilibrium depends materially on the cap, small margin and chosen response selection. It does not claim that every selection supports these prices.

\section*{Welfare and a monopolist comparison}
Use common unit weights on consumers and owners, so fees cancel. Intrinsic utility minus cost is $a-c$ for every full-coverage allocation. At the equal split, each consumer's status is $1/4$, giving status welfare
\[
                          \int_0^1\theta/4\,d\theta=1/8.
\]
An allocation concentrating all consumers with one seller gives status $1/2$ and status welfare $1/4$. It is attainable under the same technologies and prices $(B,B)$: consumers strictly prefer the populated seller for almost every positive type. Its total welfare is therefore larger by $1/8$.

A monopolist controlling the same two menus can charge $B$ and select this concentrated response, collecting total profit $B-c$. This feasible monopoly outcome exceeds the selected duopoly outcome's welfare. It does not establish a universal monopoly ranking, because the original game permits different costs, intrinsic values and multiple levels, and other duopoly selections include concentrated responses. The example establishes inefficient dominance or fragmentation only by comparing actual attainable allocations under declared welfare weights, not by using market share as a proxy.

\section*{Remaining challenge}
The existence criterion reduces seller equilibrium support to worst-response deviation values, but the response correspondence can have multiple branches and price-dependent jumps. Computing its global minima and checking all deviations for arbitrary continuous intrinsic utilities remains the main unresolved classification task. Quality choice or arbitrary mechanisms would enlarge the game further and are outside the theorem proved here.

\section{Completion Estimate}
65\%

This is a post hoc, heuristic estimate for the full catalogue target.
Consumer response existence and an exact deviation-support criterion are proved, and a credible selected duopoly equilibrium receives a concrete monopoly welfare comparison. A usable classification from arbitrary menu primitives and the complete allocation and welfare correspondence remain unresolved.

\begin{thebibliography}{9}
\bibitem{xiao} P. Xiao (2025 version), Allocating Positional Goods: A Mechanism Design Approach, arXiv:2411.06285v3. Section 2.1 and the competition discussion at the end of Section 5 inspected. \url{https://arxiv.org/html/2411.06285v3}.
\end{thebibliography}

\end{document}
